\begin{proposition}[B.4.6; Proved; a restricted non-K\"ahler cap obstruction]\label{push:B46}
Let $D$ be a compact reduced normal-crossings complex
space. Suppose its irreducible components have been
partitioned into closed unions $D_0,\ldots,D_s$, with
their reduced induced structures. Assume that
\[
 H^1(D_i,\mathbb C)\longrightarrow H^1(D_i,\mathcal O)
 \quad\hbox{is surjective for every }i,
 \tag{B.46a}
\]
and that the analogous map is injective for every
pairwise intersection $D_i\cap D_j$. Then
\[
 H^1(D,\mathbb C)\longrightarrow H^1(D,\mathcal O_D)
 \quad\hbox{is surjective}.
 \tag{B.46b}
\]

In particular these hypotheses hold for the following
specified cap assembly. The horizontal block $D_0$ is
a two-section pinching of a ruled bundle over an
analytic locally trivial torsor $V\to C$ of B.3.8,
with conductor map $A$ covering $\sigma$ and
$A^{12}=t_{6P}$. The normalization may undergo the
modifications of B.3.7, or any smooth bimeromorphic
modifications for which the two conductor restriction
maps on $H^1$ are the same pullbacks as before.
The other blocks are smooth hyperbolic toric Kato
threefolds or compact smooth spaces satisfying
(B.46a), for example class-$\mathcal C$ spaces.
Assume every pairwise overlap is either a reduced
union of invariant divisors of a hyperbolic toric
Kato threefold, or a compact normal-crossings surface
whose normalized components and normalized double
curves are rational. No such assembly can be a
semistable central fibre of the $S^6$ family.

This excludes the stated cap construction even when
an invariant Kato surface in an overlap kills a
fundamental-group generator of a Kato component.
It is not a classification of arbitrary analytic
cusp components or of their possible overlaps.
\end{proposition}
\begin{proof}
Write $D^{[p]}$ for the disjoint union of intersections
of $p+1$ distinct blocks. The augmented closed-cover
complexes resolve both $\mathbb C_D$ and $\mathcal O_D$.
For the coherent assertion this can be checked in a
normal-crossings coordinate chart. The local equation
of $D$ is a product of distinct coordinates, and each
block groups a disjoint subset of those factors.
The corresponding ideals are generated by products
in disjoint sets of variables. The ordinary alternating
intersection complex is exact, as one sees on each
monomial, where it is the augmented complex of a
simplex. The same contraction acts coefficientwise
on convergent power series. For constant sheaves,
small intersections are contractible unions of
coordinate subspaces, and the same stalkwise simplex
complex is exact. This proves both resolutions also
when the global blocks are self-crossing.

Their hypercohomology spectral sequences have
\[
 E_1^{p,q}(\mathcal F)=H^q(D^{[p]},\mathcal F),
 \qquad \mathcal F=\mathbb C\ \hbox{or}\ \mathcal O.
\]
Every connected compact reduced branch intersection
has only constant holomorphic functions. Thus the
comparison maps in the entire $q=0$ row are
isomorphisms. They induce isomorphisms on
$E_2^{1,0}$ and $E_2^{2,0}$. Surjectivity in (B.46a)
and injectivity on pairwise intersections imply
surjectivity on $E_2^{0,1}$: lift a coherent
degree-one class on the blocks, then its topological
restriction difference vanishes because its coherent
image vanishes and the pairwise comparison is
injective.

Now use the natural five-term exact sequences
\[
 0\longrightarrow E_2^{1,0}
 \longrightarrow H^1(D,\mathcal F)
 \longrightarrow E_2^{0,1}
 \xrightarrow{d_2} E_2^{2,0}.
\]
The isomorphism on the last term makes a lift of a
coherent $d_2$-closed element topologically
$d_2$-closed. The isomorphism on the first term then
gives (B.46b). If the block comparisons are
isomorphisms, the same argument also proves
injectivity, provided the map on $E_2^{0,1}$ is
injective.

We verify the horizontal-block assertion rather than
apply the pairwise criterion directly to its two
conductor branches. By B.3.8,
$H^1(V,\mathbb C)=H^1(C,\mathbb C)$ and
$H^1(V,\mathcal O)=H^1(C,\mathcal O_C)$.
The first equality is the Leray calculation there;
the second also follows from
$R^1\pi_*\mathcal O_V=\mathcal O_C(-\infty)$,
which has no global section. The action of $A$ on
these groups is $\sigma^*$, since it covers $\sigma$
and the displayed maps are pullbacks from the base.
The characteristic polynomial
\[
 p(x)=\Phi_{12}(x)\Phi_6(x),\qquad p(1)=1,
\]
shows that $1-A^*$ is invertible on both groups.
The normalization sequence of the two-section
pinching, for constants and for holomorphic
functions, therefore gives
\[
 H^1(D_0,\mathbb C)=\mathbb C,
 \qquad H^1(D_0,\mathcal O)=\mathbb C.
\]
Both classes come from the same difference map on
constant functions of the conductor. Their natural
comparison is an isomorphism. Projective-bundle
pullback identifies the normalization groups with
those of $V$. The stipulated modifications preserve
these groups and both restriction maps; in particular
the explicit modifications of B.3.7 have this property.
This grouped computation is essential: on an individual
conductor $V$, the comparison has dimensions six and
three and is not injective.

For a hyperbolic toric Kato manifold $Z$, its
fundamental group is $\mathbb Z$, as in
\cite[Section 2, p.1352]{IstratiEtAl2022Root};
\cite[Proposition 8.4 and Theorem 8.5]{IstratiEtAl2022Root}
give $H^0(Z,\Omega_Z^1)=0$ and
$H^1(Z,\mathcal O_Z)=\mathbb C$.
The exact sequence
$0\to\mathbb C\to\mathcal O\xrightarrow{d}d\mathcal O\to0$
shows that $H^1(Z,\mathbb C)\to H^1(Z,\mathcal O)$
is injective, because a global section of
$d\mathcal O$ is a closed holomorphic one-form.
The two dimensions are one, so it is an isomorphism.
The same reasoning applies to hyperbolic toric Kato
surfaces. Compact smooth rational varieties have
both groups zero. Class-$\mathcal C$ blocks satisfy
(B.46a) by their $\partial\bar\partial$ property.

Finally let $F$ be a reduced union of invariant
divisors in a hyperbolic toric Kato threefold.
Its normalized components are compact toric surfaces
or hyperbolic toric Kato surfaces
\cite[proof of Theorem 8.8, p.1380]{IstratiEtAl2022Root}.
Its normalized double curves are rational, since
every invariant curve is rational in the hyperbolic
case \cite[Theorem 5.2]{IstratiEtAl2022Root}.
Use the ordered-local-branch resolution from B.4.5.
The comparisons on component $H^1$ are isomorphisms,
and the groups on normalized double curves vanish.
The same five-term argument shows that
$H^1(F,\mathbb C)\to H^1(F,\mathcal O_F)$ is an
isomorphism. For an overlap with rational normalized
components and double curves, the argument is
identical with zero component $H^1$. Thus every
pairwise overlap in the statement has the required
injective comparison, proving (B.46b) for the
specified assembly.

For an actual semistable $S^6$ centre, B.4.2 gives
$H^1(D,\mathbb C)=0$, while upper semicontinuity
and the nearby value $h^1(\mathcal O)=1$ in B.4.3
give $h^1(D,\mathcal O_D)\geq1$. This contradicts
(B.46b). The obstruction uses the cohomology maps
on the overlaps; it does not assume that every
Kato component's loop survives in the union.
\end{proof}
